% Source-aligned English for printed p.183.
\LayerTarget{tohoku:proposition:3.10.2}
\textsc{Proposition 3.10.2.}\SourcePage{183} \textit{Suppose that the locally compact
space $X$ is countable at infinity. Then the homomorphisms \textup{(3.10.2)}, where the
projective limit is taken over the increasing filtered set of relatively compact
open subsets $U$ of $X$, are surjective. If $p>1$, for such a homomorphism to
be bijective, it is enough that for every relatively compact open $U$ there
exists another $V\supset U$ such that, for every relatively compact open $W$
containing $V$, the image of $H^{p-1}(W,F)$ in $H^{p-1}(U,F)$ under the
restriction homomorphism is identical to that of $H^{p-1}(V,F)$. If $p=1$,
for the homomorphism $H^1(X,F)\to\lim H^1(U,F)$ to be
bijective, it is enough that the groups $H^0(U,F)$ can be endowed with
complete metrizable topological-group structures for which the restriction
homomorphisms are continuous, and that for every relatively compact open $U$
there exists another $V\supset U$ such that, for every relatively compact open
$W$ containing $V$, the image of $H^0(W,F)$ in $H^0(U,F)$ is dense in that of
$H^0(V,F)$.}

\LayerTarget{tohoku:remark:proposition-3.10.2:approximation}
(Of course, in this statement one could replace relatively compact open sets
by compact sets.) This proposition, which is proved by a Mittag-Leffler
approximation procedure, is, for example, essential in the proof of the
fundamental theorems on Stein manifolds [5]. For $p=1$, it remains true if $F$
is a sheaf of groups that need not be abelian, in the following form: if the
approximation condition stated for $p=1$ holds, every element of the first
member of (3.10.2) whose image is the neutral element is itself the neutral
element.

\TohokuPartBoundary{185}
